\documentclass[10pt,a4paper]{report}

% --- Unicode fonts (compiled with XeLaTeX) ---
\usepackage{fontspec}
\setmainfont{Times New Roman}
\usepackage{hyperref}
\usepackage{graphicx}
\usepackage{tikz}
\usetikzlibrary{calc}
\usetikzlibrary{arrows.meta,positioning,fit,backgrounds}
\usepackage{array}
\usepackage{amsmath,amssymb}
\usepackage{geometry}
\usepackage{xcolor}
\usepackage{tabularx}
\usepackage{caption}
\usepackage{subcaption}
\usepackage{float}
\usepackage{enumitem}
\geometry{top=2cm,bottom=2cm,left=2.5cm,right=2.5cm}
\hypersetup{colorlinks=true,linkcolor=darkblueframe,urlcolor=darkblueframe}

\definecolor{darkblueframe}{RGB}{0,70,140}
\setcounter{secnumdepth}{3}
\setcounter{tocdepth}{2}
\setlist{itemsep=2pt,topsep=4pt}

\begin{document}
\renewcommand{\chaptername}{Chương}
\renewcommand{\contentsname}{Mục lục}
\renewcommand{\figurename}{Hình}
\renewcommand{\tablename}{Bảng}

% -------------------- Cover --------------------
\begin{titlepage}
\thispagestyle{empty}
\begin{tikzpicture}[overlay,remember picture]
  \draw[line width=1.5pt,color=darkblueframe]
    ($(current page.north west)+(1.35cm,-1.35cm)$) rectangle
    ($(current page.south east)+(-1.35cm,1.35cm)$);
  \draw[line width=0.5pt,color=darkblueframe]
    ($(current page.north west)+(1.48cm,-1.48cm)$) rectangle
    ($(current page.south east)+(-1.48cm,1.48cm)$);
\end{tikzpicture}

\centering
\vspace*{0.45cm}
{\Large\bfseries ĐẠI HỌC QUỐC GIA TP. HỒ CHÍ MINH}\\[0.18cm]
{\large\bfseries TRƯỜNG ĐẠI HỌC KHOA HỌC TỰ NHIÊN}\\[0.18cm]
{\large\bfseries KHOA CÔNG NGHỆ THÔNG TIN}

\vspace{0.35cm}
{\color{darkblueframe}\rule{0.58\textwidth}{1pt}}

\vspace{0.55cm}
\IfFileExists{logo.jpg}
  {\includegraphics[width=0.2\textwidth]{logo.jpg}}
  {\fbox{\parbox[c][2.6cm][c]{0.2\textwidth}{\centering Logo HCMUS}}}

\vspace{0.55cm}
{\color{darkblueframe}\fontsize{38}{44}\selectfont\bfseries LAB 01}

\vspace{0.25cm}
{\color{darkblueframe}\rule{0.32\textwidth}{1pt}}

\vspace{0.5cm}
{\Large\bfseries MÔ PHỎNG BLOCKCHAIN LAYER 1\\VỚI CƠ CHẾ ĐỒNG THUẬN BFT}

\vspace{0.65cm}
{\large\bfseries THÔNG TIN HỌC PHẦN}

\vspace{0.3cm}
\renewcommand{\arraystretch}{1.35}
\begin{tabular}{>{\bfseries}r @ {\hspace{0.8em}} l}
Môn học & Nhập môn blockchain và ứng dụng (CSC15010) \\
Lớp & CQ2023/22 \\
Học kỳ & Học kỳ 1 -- 2026/2027 \\
\end{tabular}

\vspace{0.55cm}
{\color{darkblueframe}\rule{0.72\textwidth}{0.6pt}}

\vspace{0.4cm}
{\large\bfseries THÀNH VIÊN THỰC HIỆN}

\vspace{0.25cm}
\begin{tabular}{>{\bfseries}r @ {\hspace{0.8em}} l}
Giảng viên & Th.S Ngô Đình Huy \\
\multicolumn{2}{c}{\rule{0pt}{0.35cm}} \\
1. & Nguyễn Gia Huy -- 23120273 \\
2. & Bùi Duy Khánh -- 231204281 \\
3. & Nguyễn Đăng Khoa -- 23120315 \\
4. & Nguyễn Ngọc Nhân Trọng -- 23120382 \\
\end{tabular}

\vfill
{\itshape TP. Hồ Chí Minh, Tháng 11 Năm 2026}
\end{titlepage}

\pagenumbering{arabic}

% -------------------- Report --------------------
\chapter{Giới thiệu}
\section{Bối cảnh}
Blockchain là một sổ cái được nhiều máy tính cùng lưu giữ. Trước khi thêm block mới, các node cần xác minh dữ liệu và thống nhất rằng block đó hợp lệ. Bài lab xây dựng một mô hình Layer 1 thu nhỏ để minh họa các thành phần chính của node validator: giao dịch có chữ ký, thực thi trạng thái, đồng thuận Byzantine Fault Tolerant (BFT), và mạng truyền thông điệp có lỗi.

Đây là simulator phục vụ học tập, không kết nối vào một blockchain công khai và không xử lý tài sản thật.

\section{Mục tiêu}
Các mục tiêu của project gồm:
\begin{itemize}
  \item Tạo và xác minh chữ ký Ed25519, băm SHA-256 và mã hóa dữ liệu nhất quán.
  \item Kiểm tra giao dịch, quyền sở hữu khóa trạng thái và chống phát lại bằng nonce.
  \item Mô hình hóa block, state root và quy trình bỏ phiếu hai pha Prevote/Precommit.
  \item Mô phỏng độ trễ, mất và gửi trùng thông điệp; tạo log có thể tái lập theo seed.
  \item Viết unit test và các kịch bản end-to-end để kiểm tra các thuộc tính an toàn cơ bản.
\end{itemize}

\section{Phạm vi}
Mã nguồn hiện thực các cấu trúc và hành vi cốt lõi ở mức mô phỏng. NetworkRunner nối transport với handler của node, timeout, proposal, hai phase vote và finalization. Mô phỏng mặc định dùng tám validator đúng; các test fault-injection kích hoạt drop, duplicate, proposer im lặng và equivocation.

\chapter{Kiến trúc hệ thống}
\section{Các tầng chính}
Project được chia thành các tầng với trách nhiệm riêng:
\begin{description}[style=nextline,leftmargin=0pt]
  \item[Crypto] \texttt{src/crypto/} quản lý khóa, chữ ký, hash và biểu diễn dữ liệu.
  \item[Execution] \texttt{src/execution/} định nghĩa transaction, block và trạng thái key-value.
  \item[Consensus] \texttt{src/consensus/} tạo phiếu, đếm quorum và giữ trạng thái khóa.
  \item[Network] \texttt{src/network/} mô phỏng message, lịch giao nhận, event loop và log sự kiện.
  \item[Node] \texttt{src/node.py} kết hợp state, consensus, block buffer và transaction pool cục bộ.
\end{description}

\begin{figure}[H]
\centering
\begin{tikzpicture}[
  box/.style={draw,rounded corners=3pt,minimum width=10.8cm,minimum height=0.72cm,align=center,font=\small},
  arr/.style={-{Stealth[length=2.2mm]},thick,color=darkblueframe},
  node distance=0.24cm
]
\node[box,fill=blue!9] (api) {Ứng dụng / RPC (lớp phục vụ được mô tả trong kiến trúc)};
\node[box,fill=orange!13,below=of api] (node) {Node validator -- \texttt{src/node.py}};
\node[box,fill=green!12,below=of node] (cons) {Consensus: Prevote, Precommit, quorum, khóa block};
\node[box,fill=yellow!18,below=of cons] (exec) {Execution: transaction, block, state root};
\node[box,fill=red!8,below=of exec] (crypto) {Crypto: chữ ký Ed25519, hash, encoding};
\node[box,fill=violet!10,below=of crypto] (net) {Network simulator: delay, drop, duplicate, logger};
\draw[arr] (api) -- (node);
\draw[arr] (node) -- (cons);
\draw[arr] (cons) -- (exec);
\draw[arr] (exec) -- (crypto);
\draw[arr,<->] (net.east) -- ++(1.1,0) |- (node.east);
\end{tikzpicture}
\caption{Sơ đồ các thành phần chính trong kiến trúc node của bài lab. Lớp RPC được nêu như vai trò kiến trúc; repo hiện chưa hiện thực module RPC riêng.}
\label{fig:node-architecture}
\end{figure}

\section{Luồng xử lý dự kiến}
Một vòng xử lý block được mô hình hóa theo chuỗi sau:
\begin{enumerate}
  \item Người dùng tạo và ký transaction bằng khóa riêng.
  \item Proposer chọn các transaction hợp lệ, thực thi thử trên bản sao state và tạo header chứa state root dự kiến.
  \item Các validator kiểm tra block, rồi phát phiếu PREVOTE cho block hash hoặc NIL.
  \item Khi có quorum PREVOTE, validator phát PRECOMMIT tương ứng.
  \item Khi đạt quorum PRECOMMIT cho cùng block, block được chốt; các transaction được áp dụng vào state.
\end{enumerate}

Với 8 validator, code tính $f=\lfloor(8-1)/3\rfloor=2$ và quorum $2f+1=5$. Đây là ngưỡng dùng trong engine để kiểm tra phiếu.

\chapter{Giao dịch, block và trạng thái}
\section{Giao dịch}
Transaction gồm sender, state key, giá trị mới, nonce và chữ ký. Payload được ký trong domain \texttt{TX:<chain\_id>}; do đó chữ ký gắn với chain cụ thể. Hàm xác minh cũng yêu cầu state key có tiền tố thuộc sender, ngăn một tài khoản sửa khóa của tài khoản khác.

Hash giao dịch được tính xác định từ payload. State từ chối transaction không hợp lệ hoặc nonce đã được dùng, nhằm giới hạn việc phát lại transaction. Các kiểm tra này được thể hiện trong \texttt{src/execution/transaction.py} và \texttt{src/execution/state.py}.

\section{Block và state root}
Header chứa parent hash, chiều cao, round, state root, transaction root và proposer. Proposer ký header theo domain \texttt{HEADER:<chain\_id>}. Block hash bắt nguồn từ payload của header. Khi tạo proposal, node thực thi thử các transaction trên bản sao state để đưa state root sau thực thi vào header.

State root là hash đại diện cho trạng thái hiện tại. Thay đổi dữ liệu hoặc nonce sẽ làm thay đổi root, giúp node so sánh trạng thái sau khi thực thi. Root hiện được tính bằng SHA-256 trên biểu diễn canonical của các cặp key--value và nonce theo sender. Repo chưa hiện thực Merkle Patricia Trie hoặc proof.

Header ký cả state root lẫn transaction root theo thứ tự transaction. Trước khi prevote, node kiểm tra chữ ký proposer, proposer theo lịch, parent hash, transaction root và state root sau khi thực thi. Khi chốt, node kiểm tra lại block và yêu cầu quorum PRECOMMIT. Nếu một transaction không hợp lệ hoặc root không khớp, node không cập nhật state và không tăng height.

\section{Lỗi phát hiện và cách xử lý}
Rà soát mã cho thấy hàm finalize cũ bỏ qua kết quả thực thi block, nên có thể tăng height dù transaction không hợp lệ. Node hiện thực thi thử trên state copy, kiểm tra root và chỉ commit sau khi có quorum PRECOMMIT. Engine cũng cần tính vote do chính validator ký vào quorum; vote cục bộ hiện được ghi nhận cùng luồng kiểm tra vote nhận. State root cũ bỏ qua nonce dù nonce ảnh hưởng transaction hợp lệ tiếp theo; root hiện cam kết cả nonce sender. Cuối cùng, header cũ không cam kết danh sách transaction; transaction root được thêm vào header và chữ ký proposer.

\chapter{Đồng thuận BFT}
\section{Hai pha bỏ phiếu}
Engine chọn proposer theo vòng luân phiên xác định từ height và round. Validator gửi PREVOTE trước; sau khi quan sát đủ phiếu cho một block, có thể gửi PRECOMMIT. Quorum được tính trên các validator duy nhất, cùng height, round và phase.

\begin{figure}[H]
\centering
\begin{tikzpicture}[
  stepbox/.style={draw,rounded corners=3pt,minimum width=3.35cm,minimum height=1.05cm,align=center,font=\small},
  arr/.style={-{Stealth[length=2.2mm]},thick,color=darkblueframe},
  node distance=0.48cm
]
\node[stepbox,fill=blue!9] (prop) {Proposer\\đề xuất block};
\node[stepbox,fill=green!12,right=of prop] (prevote) {Validators kiểm tra\\và gửi PREVOTE};
\node[stepbox,fill=yellow!18,right=of prevote] (q1) {Đạt quorum\\$2f+1$ PREVOTE};
\node[stepbox,fill=orange!15,below=0.75cm of q1] (precommit) {Gửi\\PRECOMMIT};
\node[stepbox,fill=red!9,left=of precommit] (q2) {Đạt quorum\\$2f+1$ PRECOMMIT};
\node[stepbox,fill=violet!10,left=of q2] (final) {Chốt block\\và cập nhật state};
\draw[arr] (prop) -- (prevote);
\draw[arr] (prevote) -- (q1);
\draw[arr] (q1) -- (precommit);
\draw[arr] (precommit) -- (q2);
\draw[arr] (q2) -- (final);
\end{tikzpicture}
\caption{Luồng đồng thuận hai pha được mô hình hóa trong consensus engine. Với 8 validator, $f=2$ và quorum là 5 phiếu.}
\label{fig:bft-flow}
\end{figure}

\section{Khóa và chống phiếu trùng}
ConsensusLockState ghi nhận phiếu đã ký của node và phiếu nhận từ peer. Nó từ chối phiếu từ validator không thuộc validator set, chữ ký không hợp lệ, phiếu trùng và phiếu thứ hai của cùng validator trong cùng height/round/phase. Engine chỉ tạo PRECOMMIT sau khi quan sát quorum PREVOTE; trạng thái khóa được cập nhật khi quorum PREVOTE cho một block được ghi nhận.

Nếu proposer không hoạt động, engine có thể tăng round và chọn proposer tiếp theo. Việc tăng round được mô hình hóa qua hàm chuyển round; xử lý timeout thực tế ở cấp toàn mạng còn tùy phần điều phối simulator.

\chapter{Mô phỏng mạng và ghi log}
\section{Network simulator}
NetworkSimulator dùng seed để tạo quyết định drop, delay, duplicate và thứ tự giao nhận có thể tái lập. Mỗi node có giới hạn outbound theo cửa sổ thời gian; simulator có thể block/unblock một cặp peer và log các sự kiện mạng. NetworkRunner đăng ký handler cho mỗi node và điều phối round timeout, proposal, vote và finalization.

Block được truyền theo hai bước: node nhận kiểm tra header và chữ ký trước, rồi gửi \texttt{BODY\_REQUEST}; proposer chỉ gửi body theo yêu cầu đó. Node kiểm tra lại transaction root, chữ ký transaction và state root sau khi thực thi trước khi prevote. Transaction và vote đi qua NetworkSimulator; quyết định của node chỉ dựa trên header, body và vote đã nhận tại node đó.

\section{Log xác định}
Logger ghi sự kiện theo dạng canonical để có thể so sánh byte log giữa hai lần chạy cùng seed. Khả năng tái lập giúp debug: khi kết quả khác nhau, có thể kiểm tra sự kiện nào, ở thời điểm nào, đã dẫn đến sai khác.

\chapter{Kiểm thử}
Repo tổ chức test thành \texttt{tests/unit/} và \texttt{tests/e2e/}. Unit test kiểm tra crypto, state, consensus, giới hạn outbound và peer blocking. Các scenario T1--T8 chạy qua NetworkRunner/NetworkSimulator với phạm vi sau:
\begin{description}[style=nextline,leftmargin=0pt]
  \item[T1 -- Bình thường] Proposal, header/body, hai phase vote và finalization qua mạng; mọi node chốt cùng block và state root.
  \item[T2 -- Duplicate/reorder] Bật duplicate và delay ngẫu nhiên; không đếm vote lặp và không chia finality.
  \item[T3--T4 -- Transaction lỗi/phát lại] Gửi transaction sai domain và replay qua message handler; node từ chối và log lý do.
  \item[T5 -- Drop/delay] Trước synchrony, drop/delay message; sau synchrony, các node chốt cùng block và không có hai hash final tại cùng height.
  \item[T6 -- Proposer im lặng] Bỏ qua proposer round 0; timeout đưa mạng sang round 1 và block vẫn được chốt.
  \item[T7 -- Equivocation] Hai validator ký vote mâu thuẫn gửi qua transport; mỗi node tính tối đa một vote mỗi validator/phase và không tạo quorum giả.
  \item[T8 -- Tái lập] Cùng seed tạo cùng block/state root và byte log của mọi node lẫn transport.
\end{description}

Lần chạy entry point cho kết quả 21/21 test đạt. T5 xác nhận safety và hội tụ trong cấu hình fault-injection đã chọn; đây là kiểm thử thực nghiệm hữu hạn, không phải chứng minh hình thức cho mọi lịch message.

Có thể gọi test bằng \texttt{python -m tests.run\_all\_tests}; cùng lệnh này chạy trong GitHub Actions khi push hoặc pull request. Mô phỏng một block có thể chạy bằng \texttt{python run\_simulation.py}; log canonical có thể ghi ra thư mục mới bằng tùy chọn \texttt{--write-logs}.

\chapter{Giới hạn và hướng phát triển}
\begin{itemize}
  \item Bổ sung đồng bộ/catch-up cho node bỏ lỡ nhiều height và phục hồi qua restart.
  \item Mở rộng Byzantine fault runner để tự động sinh proposal lỗi, equivocation và replay qua nhiều height.
  \item Bổ sung kiểm tra độc lập cho state root/proof và lưu bền nonce theo sender.
  \item Mở rộng ma trận fault-injection để kiểm tra nhiều seed, mức drop và cửa sổ ổn định.
\end{itemize}

\chapter{Kết luận}
Project cung cấp mô hình cơ bản để học cách một node blockchain xử lý transaction, tạo block, bỏ phiếu BFT và cập nhật state. Việc tách crypto, execution, consensus và network giúp từng mối quan tâm dễ đọc và kiểm tra hơn. Các test hiện có minh họa nhiều tình huống quan trọng như chữ ký sai, transaction bị phát lại, phiếu trùng và log xác định.

NetworkRunner hiện điều phối luồng một height qua mạng mô phỏng. Bước tiếp theo là bổ sung catch-up khi node bỏ lỡ height và mở rộng fault-injection nhiều height, nhiều lịch message.

\section*{Tài liệu tham khảo}
\begin{enumerate}[leftmargin=*,itemsep=0pt,topsep=2pt]
  \sloppy
  \item Blockchain course, \emph{Lab 01} (đề bài cung cấp, 2026).
  \item Knowledge Engineering Department, HCMUS, \emph{Lab 01: Node Architecture, Clients, and State Synchronization}, slides, 6 Nov 2025.
  \item cryptography contributors, \emph{Ed25519 signing}, \url{https://cryptography.io/en/latest/hazmat/primitives/asymmetric/ed25519/}.
  \item pytest contributors, \emph{pytest documentation}, \url{https://docs.pytest.org/en/stable/}.
  \item GitHub, \emph{Workflow syntax for GitHub Actions}, \url{https://docs.github.com/en/actions/writing-workflows/workflow-syntax-for-github-actions}.
\end{enumerate}

\end{document}
